\section{Discussion}

Our results validate temporal ordering---spurious features converge 4 epochs earlier than core features on CMNIST via gradient-level measurement---while identifying cross-dataset $\rho_j$ transfer as a fundamental constraint for neuron-level interventions. We interpret findings, acknowledge limitations, and discuss broader implications.

\subsection{Key Findings Interpretation}

\textbf{Temporal ordering is measurable and mechanistically grounded.} Spurious-only training converges at $E_s=13$, core-only at $E_c=17$, yielding $\Delta=4$ epochs temporal gap (h-e1). This substantially exceeds predicted threshold ($\Delta \geq 2$) in proof-of-concept validation, providing first direct gradient-level validation of JTT/LfF hypothesis. Layer-wise neuron-spurious correlation analysis (h-m1) confirms mechanism: early layers show significantly higher $\rho_j$ than late layers ($p=0.0028$, Cohen's $d=0.25$), consistent with simpler low-level features (color) being processed earlier in CNN hierarchy than high-level semantic features (shape). Temporal gap is not dataset artifact but architectural phenomenon amplified by hierarchical layer-wise processing.

\textbf{Reweighting methods succeed via temporal misalignment correction.} JTT trains biased ERM model, identifies hard examples (those ERM misclassifies), then retrains with hard examples upweighted. Our temporal gap measurement explains why this works: hard examples are those requiring core features (learned late, after epoch 13), while easy examples exploit spurious features (learned early, by epoch 13). Upweighting late-learned examples implicitly targets core-feature-dependent samples, correcting for temporal misalignment. This mechanistic evidence connects optimization dynamics (implicit bias toward simpler features) to empirical debiasing success (reweighting late-learned examples).

\textbf{Neuron-level interventions face dataset-specificity constraint.} Gradient-aware training (h-c1) using CMNIST $\rho_j$ on Waterbirds failed catastrophically (39\% vs 86\% JTT target), revealing that $\rho_j$ values are dataset-specific and spurious-feature-type dependent. Color spurious correlation (CMNIST) produces different neuron activation patterns than background spurious correlation (Waterbirds). Network learns dataset-specific feature representations, so $\rho_j$ computed on one dataset cannot be reused on another. This negative result identifies theoretical boundary: gradient-aware methods require per-dataset calibration phase (re-run h-m1 on target dataset), limiting scalability compared to example-level reweighting (JTT).

\subsection{Honest Limitations}

\textbf{Single-dataset validation (CMNIST only).} Temporal ordering validated on CMNIST color spurious correlation; generalization to background (Waterbirds), attribute (CelebA), context (NICO++) spurious types remains untested (75\% scope reduction from original 4-dataset plan). CMNIST is canonical spurious benchmark (Arjovsky et al. 2019), but ablation methodology (blur removes shape, grayscale removes color) is CMNIST-specific. Waterbirds requires background segmentation masks, CelebA requires gender-balanced sampling, NICO++ requires context masking---each dataset needs custom ablation strategy. Manual setup barrier (approximately 2-4 hours per dataset) deferred multi-dataset validation to future work. \textbf{Why acceptable:} CMNIST proof-of-concept establishes mechanism existence before resource-intensive multi-dataset sweep. Ablation training principle (isolate feature types, measure convergence) is dataset-agnostic, extension straightforward. \textbf{Mitigation:} Multi-dataset validation (FW-C1) planned, requires manual dataset preprocessing but uses identical convergence measurement protocol (gradient norm $<$ 10\% peak for 3 epochs).

\textbf{Proof-of-concept statistical validation pending.} h-e1 reported results use seed 0 only; full 10-seed validation launched but not completed by submission deadline. Temporal gap $\Delta=4$ epochs observed on single seed, substantially exceeding threshold ($\Delta \geq 2$), provides directional evidence but p-value and confidence intervals not established. Large effect size (4 vs 2 epochs) suggests seed variance unlikely to eliminate gap, but statistical significance unconfirmed. \textbf{Why acceptable:} MUST\_WORK gate requires proof-of-concept, not full statistical validation. $\Delta=4$ margin provides buffer against seed variance---even if worst seed shows $\Delta=2.5$ epochs, threshold still met. \textbf{Mitigation:} Full 10-seed validation (FW-B1) in progress (approximately 60 hours compute), will establish mean $\Delta$, std dev, paired t-test p-value.

\textbf{GradCAM diagnostic boundary.} GradCAM temporal ratio (h-e3) failed---$R_{\text{temporal}}$ increased (0.478 $\rightarrow$ 0.489) instead of decreasing as predicted. Root cause: spatial masking incompatibility with global corruption. CMNIST color applied uniformly across entire image (not spatially localized), so center/outer spatial masks do not separate color from shape. h-e3 designed for Waterbirds (spatially separated background vs foreground bird), wrong dataset choice. \textbf{Why acceptable:} Identifies clear scope boundary---$R_{\text{temporal}}$ requires spatial separation between spurious and core feature regions. Does not refute temporal ordering (h-e1 validated via ablation, independent of GradCAM). Method-specific failure, not hypothesis failure. \textbf{Mitigation:} Retest h-e3 on Waterbirds with proper background segmentation masks (FW-A1). Consider alternative attribution methods (Integrated Gradients, SHAP) less dependent on spatial assumptions.

\textbf{Gradient variance metric artifact.} Variance ratio test (h-e2) failed: $V_{\text{spurious}} / V_{\text{core}} = 0.77 > 0.7$ threshold (10\% miss). Post-convergence zero variance in spurious-only variant (converged epoch 10, measured through epoch 30) skewed ratio calculation upward---spurious variant stopped training early, zero gradients afterward inflated mean variance ratio. \textbf{Why acceptable:} DESIGN\_ISSUE not HYPOTHESIS\_ISSUE. Alternative metric formulation (variance during active training window only, excluding post-convergence) may validate claim. Forgetting rate already provides stability evidence ($F_{\text{spurious}} < F_{\text{core}}$), variance is supplementary. \textbf{Mitigation:} Recompute variance using only active epochs (spurious: 1-10, core: 1-17) or use gradient norm decay rate instead of variance (FW-B2).

\subsection{Broader Impact}

\textbf{For research community:} Mechanistic evidence explaining why JTT/LfF succeed provides foundation for next-generation debiasing methods. Adaptive reweighting schedules could monitor gradient convergence in real-time, adjusting example weights dynamically as temporal gap emerges (rather than fixed $T_{\text{up}}$ upweighting epoch in JTT). Architectural interventions might enforce late-layer learning delays (e.g., freeze early layers initially, unfreeze after epoch 10) to prevent early spurious dominance.

\textbf{For practitioners:} Temporal ordering validated, but neuron-level modulation (gradient-aware training) faces dataset-specificity constraints. Simpler example-level reweighting (JTT) remains more practical---no per-dataset calibration required, transfers across spurious types. Gradient-aware methods viable within-dataset (e.g., continual learning on same domain with drift), but cross-dataset transfer requires re-running h-m1 analysis (additional compute burden).

\textbf{Fairness implications:} Medical diagnosis, lending, hiring deploy models where spurious correlations (patient demographics correlating with diagnosis, ZIP code correlating with creditworthiness) cause disparate performance across groups. Validated temporal ordering enables practitioners to monitor training dynamics---if worst-group accuracy plateaus early while average accuracy improves, temporal gap emerging. Early detection allows intervention (switch to reweighting schedule) before convergence locks in spurious solution.
